\chapter{粒子、場與相對論記號}

\section{課前的多普勒頻移問題}

討論從上週留下的一個問題開始：以接近光速運動的觀察者，測得的波長是否總是較長？不是，答案取決於運動方向。觀察者若與光波同向運動，會看到波長變長；若逆著光波運動，則會看到波長變短。這就是相對論性多普勒效應。

\begin{classroomqa}
\audiencequestion{接近光速運動會使觀測到的波長變長嗎？}
\lecturerresponse{只有在觀察者與光同向運動時才會如此。反轉觀察者的速度方向，頻移方向也會反轉。這個結果可以直接由洛倫茲幾何推得，不必另外背誦一條公式。\lecturetimestamp{00:00:13}}
\end{classroomqa}

採用\(c=1\)的單位制，並令帶撇號的參考系以速率\(v\)向右運動。向右傳播的光波波峰沿著零線前進。將未加撇號參考系中的波長歸一化為 1，並把兩個相鄰波峰的世界線寫成
\begin{equation}
x=t,
\qquad
x=t+1.
\label{eq:sr5-crest-lines}
\end{equation}
在未加撇號的參考系中，要在相同的\(t\)下測量兩個波峰之間的距離。帶撇號參考系中的波長則必須在相同的\(t'\)下測量。因此，整個計算的關鍵就在於同時性的相對性。

\begin{figure}[htbp]
\centering
\includegraphics[width=0.90\textwidth]{lecture_05_figure_01.png}
\caption{相鄰的右行波峰、光錐，以及運動參考系傾斜的座標軸。}
\lectureframe{00:03:30}
\end{figure}

推導在此暫停，先簡短討論課程安排。暑期不會開課。下一門課可能講廣義相對論，也可能更全面地回到量子力學；特別提到，在學習量子場論之前，應先理解諧振子。

\begin{classroomqa}
\audiencequestion{未來的課程能否討論紅外--紫外關聯？}
\lecturerresponse{這個問題涉及量子力學與引力的交會。本課程所學的量子力學已足以讓大家認識這個問題，但必須先建立引力理論；最自然的途徑是學習廣義相對論。\lecturetimestamp{00:05:37}}
\end{classroomqa}

\begin{classroomqa}
\audiencequestion{能否直接從黑板上的圖讀出傾斜座標軸的相對尺度？}
\lecturerresponse{不行。時空圖呈現的是因果關係與同時性的幾何，但用一般的歐幾里得直尺量圖上的距離，並不等於進行洛倫茲測量。應先找出相關事件，再套用洛倫茲變換。\lecturetimestamp{00:06:10}}
\end{classroomqa}

現在回到計算。將未加撇號參考系中的波峰間距歸一化為 1，再用代數求出帶撇號參考系中同時發生的兩個事件。

洛倫茲變換為
\begin{equation}
t'=\gamma(t-vx),
\qquad
x'=\gamma(x-vt),
\qquad
\gamma=\frac{1}{\sqrt{1-v^2}}.
\label{eq:sr5-lorentz-doppler}
\end{equation}
選取帶撇號參考系的同時性切片\(t'=0\)。由方程
\eqref{eq:sr5-lorentz-doppler}可得
\begin{equation}
t=vx.
\label{eq:sr5-primed-simultaneity}
\end{equation}
第一個波峰在原點與此切片相交。第二個波峰滿足
\(x=t+1\)，因此其交點滿足
\begin{align}
x&=vx+1,\\
x&=\frac{1}{1-v},
\qquad
t=\frac{v}{1-v}.
\end{align}
其帶撇號座標為
\begin{align}
x'
&=\frac{x-vt}{\sqrt{1-v^2}}\\
&=\frac{1-v^2}{(1-v)\sqrt{1-v^2}}\\
&=\sqrt{\frac{1+v}{1-v}}.
\end{align}
原來的波長為 1，因此
\begin{equation}
\boxed{
\frac{\lambda'}{\lambda}
=
\sqrt{\frac{1+\beta}{1-\beta}}
},
\qquad
\beta=\frac{v}{c}.
\label{eq:sr5-doppler}
\end{equation}
等價地，
\begin{equation}
\frac{\lambda'}{\lambda}
=
\sqrt{\frac{c+v}{c-v}}.
\end{equation}

\begin{figure}[htbp]
\centering
\includegraphics[width=0.90\textwidth]{lecture_05_figure_02.png}
\caption{恢復\(c\)後的多普勒因子；反轉\(v\)，頻移因子便變為其倒數。}
\lectureframe{00:13:40}
\end{figure}

當\(v>0\)時，帶撇號參考系中的觀察者追著光前進，測得較長的波長。將\(v\)換成\(-v\)，就表示觀察者逆著光運動，此時頻移因子小於 1。如果觀察者追著光線並能接近\(c\)，波長就會發散。頻率的變化比則是波長比的倒數。

\section{為什麼相互作用必須出現在同一個作用量中}

前面的問題告一段落後，主題轉向場與粒子。在量子場論中，粒子與場最終會成為同一類激發的兩種描述，但此處採用的並非這種架構。在這裡，粒子沿著世界線運動，經典場則分布於時空中。眼前的問題較為簡單：如果場作用於粒子，粒子又如何反過來作用於場？

作用量原理透過一個有限維的原型回答這個問題。令\(x(t)\)和\(y(t)\)為兩個廣義座標。若拉格朗日量可以分離，
\begin{equation}
L(x,\dot x,y,\dot y)
=
L_x(x,\dot x)+L_y(y,\dot y),
\label{eq:sr5-separable-l}
\end{equation}
那麼兩個尤拉--拉格朗日方程便彼此獨立：
\begin{equation}
\frac{d}{dt}\frac{\partial L}{\partial\dot x}
=
\frac{\partial L}{\partial x},
\qquad
\frac{d}{dt}\frac{\partial L}{\partial\dot y}
=
\frac{\partial L}{\partial y}.
\end{equation}
\(x\)的方程不含\(y\)，而\(y\)的方程不含\(x\)。

若要讓這些變數彼此作用，就加入一個不能拆成\(x\)部分與\(y\)部分的項。簡單的例子是
\begin{equation}
L
=
\frac{\dot x^2}{2}
+\frac{\dot y^2}{2}
-V(x)-W(y)+\kappa xy.
\label{eq:sr5-coupled-coordinates}
\end{equation}
現在，對\(x\)變分會產生一個與\(y\)成正比的項；對\(y\)變分則會產生一個與\(x\)成正比的項。同一個混合項決定了兩個方向的作用。

\begin{figure}[htbp]
\centering
\includegraphics[width=0.90\textwidth]{lecture_05_figure_03.png}
\caption{拉格朗日量中的一個項若同時包含兩個廣義座標，兩者便會產生動力學耦合。}
\lectureframe{00:19:00}
\end{figure}

這正是此處所說的作用--反作用關係：兩個物件都是動力學變數，並且都在同一個作用量中接受變分。若將場視為外加的背景，場的響應可能就不會顯現；但場一旦本身也是動力學變數，耦合項便會同時導出兩個方程。

\section{標量場如何作用於粒子}

\subsection{相對論性耦合}

先把標量場\(\phi(\mathbf{x},t)\)視為已知。它可以是波，也可以是靜態分布；目前暫不考慮它自身的方程。自由相對論性粒子的作用量為
\begin{equation}
S_{\mathrm p}
=
-m\int d\tau.
\end{equation}
若粒子沿直線運動，且\(c=1\)，則
\begin{equation}
L_{\mathrm p}^{(0)}
=
-m\sqrt{1-\dot x^2}.
\end{equation}
最簡單的標量耦合，是以依賴場值的係數取代質量係數：
\begin{equation}
\boxed{
L_{\mathrm p}
=
-\bigl[m+g\phi(x,t)\bigr]\sqrt{1-\dot x^2}
}.
\label{eq:sr5-particle-coupling}
\end{equation}
耦合常數\(g\)表示場對此粒子的作用強度。其正負號是一種約定：將\(g\)改為\(-g\)，標量力的方向便會反轉。

\begin{figure}[htbp]
\centering
\includegraphics[width=0.90\textwidth]{lecture_05_figure_04.png}
\caption{將\(m\)換成局域係數\(m+g\phi(x,t)\)後的自由粒子作用量。}
\lectureframe{00:24:10}
\end{figure}

\subsection{慢速粒子極限}

恢復光速因子，讓牛頓極限清楚顯現：
\begin{equation}
L_{\mathrm p}
=
-\bigl[mc^2+g\phi(x,t)\bigr]
\sqrt{1-\frac{\dot x^2}{c^2}}.
\label{eq:sr5-particle-coupling-c}
\end{equation}
當\(|\dot x|\ll c\)時，
\begin{equation}
\sqrt{1-\frac{\dot x^2}{c^2}}
=
1-\frac{\dot x^2}{2c^2}
+O\!\left(\frac{\dot x^4}{c^4}\right).
\end{equation}
將此式代入\eqref{eq:sr5-particle-coupling-c}，得到
\begin{align}
L_{\mathrm p}
&=
-mc^2-g\phi
+\frac{m\dot x^2}{2}
+\frac{g\phi\,\dot x^2}{2c^2}
+O(c^{-4})\\
&\longrightarrow
\frac{m\dot x^2}{2}-g\phi(x,t),
\label{eq:sr5-nr-particle}
\end{align}
略去常數\(-mc^2\)以及受\(c^{-2}\)抑制的項後，結果就是熟悉的形式\(T-V\)，因此
\begin{equation}
V_{\mathrm p}(x,t)=g\phi(x,t).
\label{eq:sr5-scalar-potential}
\end{equation}

\begin{figure}[htbp]
\centering
\includegraphics[width=0.90\textwidth]{lecture_05_figure_05.png}
\caption{平方根展開後，留下牛頓力學的動能項與標量勢能\(g\phi\)。}
\lectureframe{00:28:50}
\end{figure}

龐大的靜止能量常數不會進入尤拉--拉格朗日方程，但依賴場值的項會進入。在電磁學中，電荷扮演類似的耦合強度角色；不過，實際的電磁相互作用使用的是四維向量勢，而不是這個標量玩具模型。

\section{粒子如何反過來作用於場}

\subsection{一個變分問題}

現在不能再把場單純視為已知背景。考慮一個時空區域，其中同時包含一個場位形與一條粒子世界線。物理解會使總作用量在對兩者變分時取駐值：
\begin{equation}
\delta S[\phi,x]=0.
\end{equation}
固定場而對世界線變分，會得到粒子方程；固定世界線而對場變分，則會得到場方程。連結兩者的，正是它們共同的相互作用項。

\begin{figure}[htbp]
\centering
\includegraphics[width=0.90\textwidth]{lecture_05_figure_06.png}
\caption{粒子世界線穿過標量場所在的時空區域。}
\lectureframe{00:30:30}
\end{figure}

自由無質量標量場的作用量取為
\begin{equation}
S_{\mathrm f}
=
\int d^4x\,
\frac12
\left[
(\partial_t\phi)^2
-(\partial_x\phi)^2
-(\partial_y\phi)^2
-(\partial_z\phi)^2
\right].
\label{eq:sr5-free-field-action}
\end{equation}
從這裡開始，再次採用\(c=1\)。若寫出明確的單位，時間項會帶有相應的\(1/c^2\)因子，實際因子取決於座標與場的歸一化方式。

\begin{figure}[htbp]
\centering
\includegraphics[width=0.90\textwidth]{lecture_05_figure_07.png}
\caption{場作用量涵蓋的時空區域與粒子世界線所在區域相同。}
\lectureframe{00:32:10}
\end{figure}

\subsection{靜止粒子與世界線積分}

為了用最少的代數運算看出源項，令粒子靜止於原點。此時因子\(d\tau\)等於\(dt\)，而\eqref{eq:sr5-particle-coupling}中的相互作用部分變為
\begin{equation}
S_{\mathrm{int}}
=
-g\int dt\,\phi(\mathbf{0},t).
\label{eq:sr5-worldline-interaction}
\end{equation}
這是沿世界線的積分；相較之下，\eqref{eq:sr5-free-field-action}是對整個時空積分。若要將兩者合併為一個時空積分，就需要一個能選出粒子所在位置的函數。

\begin{figure}[htbp]
\centering
\includegraphics[width=0.90\textwidth]{lecture_05_figure_08.png}
\caption{場區域中的靜止世界線，以及要與之合併的場作用量。}
\lectureframe{00:35:35}
\end{figure}

\subsection{狄拉克 delta 函數}

一維 delta 函數定義如下：
\begin{equation}
\int_{-\infty}^{\infty}dx\,f(x)\delta(x-a)=f(a).
\label{eq:sr5-delta-definition}
\end{equation}
它在\(x=a\)以外為零，在\(x=a\)處是奇異的，且積分為 1。最好將它理解為一種分布，其特性由\eqref{eq:sr5-delta-definition}刻畫，而不是一般的有限值函數。

\begin{figure}[htbp]
\centering
\includegraphics[width=0.90\textwidth]{lecture_05_figure_09.png}
\caption{將原點處的場值改寫為空間積分的 delta 函數恆等式。}
\lectureframe{00:36:45}
\end{figure}

在三維空間中，
\begin{equation}
\delta^3(\mathbf{x})
=
\delta(x)\delta(y)\delta(z),
\end{equation}
因此
\begin{align}
\int d^3x\,
\phi(\mathbf{x},t)\delta^3(\mathbf{x})
&=
\int dx\,dy\,dz\,
\phi(x,y,z,t)\\
&\quad\times\delta(x)\delta(y)\delta(z)\\
&=\phi(\mathbf{0},t).
\label{eq:sr5-triple-delta}
\end{align}

\begin{classroomqa}
\audiencequestion{\(d^3x\)是從哪裡來的？}
\lecturerresponse{這就是一般的空間體積元素\(dx\,dy\,dz\)。在一維空間中，同一個恆等式只包含\(dx\)和一個 delta 函數。\lecturetimestamp{00:37:09}}
\end{classroomqa}

\begin{figure}[htbp]
\centering
\includegraphics[width=0.90\textwidth]{lecture_05_figure_10.png}
\caption{三維 delta 函數展開為三個一維因子的乘積。}
\lectureframe{00:39:05}
\end{figure}

現在可以將方程\eqref{eq:sr5-worldline-interaction}寫成
\begin{equation}
S_{\mathrm{int}}
=
-g\int d^4x\,
\delta^3(\mathbf{x})\phi(\mathbf{x},t).
\end{equation}
完整的靜態源作用量為
\begin{align}
S
=
\int d^4x
\biggl\{
&\frac12(\partial_t\phi)^2
-\frac12|\boldsymbol{\nabla}\phi|^2
\nonumber\\
&-g\delta^3(\mathbf{x})\phi
\biggr\}.
\label{eq:sr5-combined-action}
\end{align}

\begin{classroomqa}
\audiencequestion{將相互作用項納入時空作用量時，是否漏掉了負號？}
\lecturerresponse{世界線拉格朗日量中有\(-g\phi\)，因此時空作用密度中也必須有\(-g\delta^3(\mathbf{x})\phi\)。保留這個負號，才能得到場方程中的源項。\lecturetimestamp{00:41:24}}
\end{classroomqa}

\begin{figure}[htbp]
\centering
\includegraphics[width=0.90\textwidth]{lecture_05_figure_11.png}
\caption{自由標量場的作用密度與局域相互作用項合併為同一個時空作用量。}
\lectureframe{00:42:05}
\end{figure}

\subsection{有源場方程}

對\eqref{eq:sr5-combined-action}關於\(\phi\)變分，可得
\begin{equation}
\boxed{
\partial_t^2\phi-\nabla^2\phi
=
-g\delta^3(\mathbf{x})
}.
\label{eq:sr5-sourced-wave}
\end{equation}
在原點以外，這就是自由波動方程；原點處的 delta 函數則標示出粒子這個局域源。

\begin{figure}[htbp]
\centering
\includegraphics[width=0.90\textwidth]{lecture_05_figure_12.png}
\caption{由原點處的點源驅動的波動算符。}
\lectureframe{00:45:35}
\end{figure}

靜態解滿足\(\partial_t\phi=0\)，所以
\begin{equation}
\boxed{
\nabla^2\phi
=
g\delta^3(\mathbf{x})
}.
\label{eq:sr5-poisson}
\end{equation}
這是點源的泊松方程。此處的類比是靜電學，還不是電磁學：這個場仍是標量\(\phi\)，而非電磁四維勢。

\begin{figure}[htbp]
\centering
\includegraphics[width=0.90\textwidth]{lecture_05_figure_13.png}
\caption{有源波動方程的靜態極限，即點源的泊松方程。}
\lectureframe{00:47:05}
\end{figure}

互易關係現在已清楚呈現：項\(g\phi\)給粒子勢能，而同一個\(g\)又使粒子成為\(\phi\)的源。這兩種作用並非彼此獨立的假設。

\subsection{運動源會輻射}

固定在\(\mathbf{a}\)的粒子以\(\delta^3(\mathbf{x}-\mathbf{a})\)表示。若其位置隨時間改變，源項就變為
\begin{equation}
\delta^3\!\left(\mathbf{x}-\mathbf{a}(t)\right).
\label{eq:sr5-moving-delta}
\end{equation}
含時的右側無法由不隨時間變化的場來平衡，因此\(\partial_t^2\phi\)項必須重新出現。加速或振盪的源可以在場中激發波。

\begin{figure}[htbp]
\centering
\includegraphics[width=0.90\textwidth]{lecture_05_figure_14.png}
\caption{delta 函數源的位置移至\(\mathbf{a}(t)\)，使場方程具有時間相依性。}
\lectureframe{00:50:30}
\end{figure}

\begin{classroomqa}
\audiencequestion{若將源視為非相對論性的，場方程中的相對論性又在哪裡？}
\lecturerresponse{相對論性的傳播由左側的時空波動算符體現。完整的相對論性粒子描述會保留固有時作用量與速度因子。即使源運動緩慢，只要隨時間變化，也能激發波。\lecturetimestamp{00:52:14}}
\end{classroomqa}

\begin{classroomqa}
\audiencequestion{\(g\phi\)的單位是什麼？}
\lecturerresponse{在粒子作用量中，這一項與\(mc^2\)並列，因此具有能量單位。在此採用的自然單位制歸一化下，\(3+1\)維時空中的標量場，其質量量綱為 1，而\(g\)可以是無量綱的。量綱會隨時空維度與場的歸一化方式而改變。\lecturetimestamp{00:52:55}}
\end{classroomqa}

\section{四維向量、標量與度規}

分量方程已經變得相當冗長，因此值得改用更簡潔的記號。將時空座標寫成
\begin{equation}
x^\mu=(x^0,x^1,x^2,x^3)=(t,x,y,z),
\label{eq:sr5-four-coordinate}
\end{equation}
其中希臘指標涵蓋全部四個分量。四個分量組成的列表不一定就是四維向量；只有當其分量依照\(x^\mu\)在洛倫茲變換下的方式變換時，它才是四維向量。

\begin{figure}[htbp]
\centering
\includegraphics[width=0.90\textwidth]{lecture_05_figure_15.png}
\caption{以\(x^\mu\)表示的四個時空座標。}
\lectureframe{00:55:00}
\end{figure}

三個空間分量為零，可能只在某一個參考系中成立，在另一個參考系中則不成立。相對地，若四個分量全為零，所有洛倫茲觀察者都會同意該向量為零。像\(dx^\mu\)這樣的差值也是四維向量。

不變的固有時區間為
\begin{equation}
d\tau^2
=
dt^2-dx^2-dy^2-dz^2.
\label{eq:sr5-proper-time}
\end{equation}
它只有一個值，並非由四個依參考系而異的分量構成，因此是標量。

\begin{figure}[htbp]
\centering
\includegraphics[width=0.90\textwidth]{lecture_05_figure_16.png}
\caption{四維位移與不變的固有時區間。}
\lectureframe{00:57:25}
\end{figure}

\subsection{閔可夫斯基度規}

採用時間分量為負、空間分量為正的號差慣例：
\begin{equation}
\eta_{\mu\nu}
=
\operatorname{diag}(-1,1,1,1).
\label{eq:sr5-metric}
\end{equation}
時間分量為正、空間分量為負的另一種慣例也同樣有效，但每個方程都必須前後一致。根據\eqref{eq:sr5-metric}，
\begin{equation}
dx^\mu dx_\mu
=
-dt^2+dx^2+dy^2+dz^2
=
-d\tau^2.
\end{equation}
整體符號不會影響不變性。

\begin{figure}[htbp]
\centering
\includegraphics[width=0.90\textwidth]{lecture_05_figure_17.png}
\caption{採用時間分量為負、空間分量為正號差的閔可夫斯基度規，以對角矩陣表示。}
\lectureframe{01:00:30}
\end{figure}

度規可以降低指標：
\begin{equation}
A_\mu=\eta_{\mu\nu}A^\nu.
\label{eq:sr5-lower-index}
\end{equation}
對這個笛卡兒座標系中的度規而言，
\begin{equation}
A^\mu=(A^0,A^1,A^2,A^3),
\qquad
A_\mu=(-A^0,A^1,A^2,A^3).
\label{eq:sr5-upper-lower-components}
\end{equation}
再套用一次相同的度規就會升高指標，因為\(\eta^{\mu\rho}\eta_{\rho\nu}=\delta^\mu{}_\nu\)。

\begin{figure}[htbp]
\centering
\includegraphics[width=0.90\textwidth]{lecture_05_figure_18.png}
\caption{以矩陣乘法乘上\(\eta_{\mu\nu}\)來降低指標。}
\lectureframe{01:03:00}
\end{figure}

\begin{classroomqa}
\audiencequestion{這裡的矩陣乘法真的包含求和嗎？}
\lecturerresponse{是的。每個輸出分量都是\(\eta\)的一列與整個向量相乘，也就是對欄指標求和。度規是對角矩陣，因此計算看起來很簡單；一般矩陣則會明確顯示這個求和。\lecturetimestamp{01:04:56}}
\end{classroomqa}

\subsection{愛因斯坦求和約定與自由指標}

愛因斯坦求和約定省略求和符號：只要同一個指標在上標與下標各出現一次，就表示要對該指標求和：
\begin{equation}
\eta_{\mu\nu}A^\nu
\equiv
\sum_{\nu=0}^{3}\eta_{\mu\nu}A^\nu.
\end{equation}
這裡的\(\nu\)是虛指標，求和後便消去，也可以改用其他名稱。\(\mu\)是自由指標，標示結果的四個分量。一個指標通常不應在同一個單項式中出現超過兩次。

傳統上，上標指標稱為\emph{逆變}，下標指標稱為\emph{協變}。本講刻意採用較直白的說法：\emph{上}標與\emph{下}標；不過，兩套術語指的是相同的變換規則。

\section{縮並與標量的導數}

\subsection{閔可夫斯基內積}

將向量與其降指標後的形式縮並，得到
\begin{equation}
A^\mu A_\mu
=
\eta_{\mu\nu}A^\mu A^\nu
=
-(A^0)^2+(A^1)^2+(A^2)^2+(A^3)^2.
\label{eq:sr5-vector-square}
\end{equation}
不必另外手動加上負號；它已經包含在降指標後的時間分量中。

\begin{figure}[htbp]
\centering
\includegraphics[width=0.90\textwidth]{lecture_05_figure_19.png}
\caption{縮並\(A^\mu A_\mu\)展開為時間分量與空間分量。}
\lectureframe{01:09:30}
\end{figure}

\begin{classroomqa}
\audiencequestion{下標向量應該畫成列向量，而不是行向量嗎？}
\lecturerresponse{在初等矩陣記號中，可以把它表示成列向量，但指標記號並不要求採用這種表示法。關鍵在於它的變換方式，以及它如何與上標向量縮並。\lecturetimestamp{01:10:26}}
\end{classroomqa}

\begin{classroomqa}
\audiencequestion{希臘指標總是涵蓋全部四個分量嗎？}
\lecturerresponse{在這套記號中，是的。\(\mu,\nu\)等希臘指標涵蓋時空分量；\(i,j\)等拉丁指標則通常只涵蓋三個空間分量。\lecturetimestamp{01:12:23}}
\end{classroomqa}

同樣的構造也適用於兩個不同的四維向量：
\begin{equation}
A^\mu B_\mu
\end{equation}
這是一個標量。簡短的證明可以使用極化恆等式。由於\((A+B)^2\)與\((A-B)^2\)都是標量，
\begin{equation}
(A+B)^2-(A-B)^2=4A^\mu B_\mu
\end{equation}
也必然是標量。這就是普通內積的閔可夫斯基版本。

\begin{classroomqa}
\audiencequestion{\(A^\mu B_\mu\)類似普通的內積嗎？}
\lecturerresponse{沒錯，只是閔可夫斯基度規在時間分量與空間分量之間引入了相對符號。\lecturetimestamp{01:15:44}}
\end{classroomqa}

\subsection{梯度是協向量}

令\(\phi(x)\)為標量場，其中\(x\)現在代表全部四個時空座標。在相鄰事件之間，
\begin{equation}
d\phi
=
\frac{\partial\phi}{\partial x^\mu}dx^\mu
=
\partial_\mu\phi\,dx^\mu.
\label{eq:sr5-dphi}
\end{equation}
左側是標量，而\(dx^\mu\)是四維向量。因此，係數
\begin{equation}
\partial_\mu\phi
=
\left(
\partial_t\phi,
\partial_x\phi,
\partial_y\phi,
\partial_z\phi
\right)
\label{eq:sr5-gradient-lower}
\end{equation}
必須依協向量的方式變換。升高其指標可得
\begin{equation}
\partial^\mu\phi
=
\left(
-\partial_t\phi,
\partial_x\phi,
\partial_y\phi,
\partial_z\phi
\right).
\label{eq:sr5-gradient-upper}
\end{equation}

\begin{figure}[htbp]
\centering
\includegraphics[width=0.90\textwidth]{lecture_05_figure_20.png}
\caption{微分\(d\phi=\partial_\mu\phi\,dx^\mu\)，由此可看出標量梯度具有協向量的性質。}
\lectureframe{01:20:00}
\end{figure}

\begin{classroomqa}
\audiencequestion{\(x^\mu\)中的上標，出現在分母\(\partial/\partial x^\mu\)時會變成下標嗎？}
\lecturerresponse{會。座標微分\(dx^\mu\)依逆變方式變換，因此其對偶導數\(\partial_\mu=\partial/\partial x^\mu\)依協變方式變換。這使得\(d\phi=\partial_\mu\phi\,dx^\mu\)成為標量。\lecturetimestamp{01:20:00}}
\end{classroomqa}

\subsection{不變的標量場作用密度}

將兩個導數向量縮並：
\begin{align}
\partial_\mu\phi\,\partial^\mu\phi
&=
-(\partial_t\phi)^2
+(\partial_x\phi)^2
+(\partial_y\phi)^2
+(\partial_z\phi)^2.
\end{align}
因此，自由標量場的作用密度可以簡潔地寫成
\begin{equation}
\boxed{
\mathcal L_0
=
-\frac12
\partial_\mu\phi\,\partial^\mu\phi
}.
\label{eq:sr5-compact-density}
\end{equation}
這與\eqref{eq:sr5-free-field-action}中的分量形式完全相同。

\begin{figure}[htbp]
\centering
\includegraphics[width=0.90\textwidth]{lecture_05_figure_21.png}
\caption{標量場的動能密度可濃縮為一個縮並表示式。}
\lectureframe{01:25:10}
\end{figure}

\begin{classroomqa}
\audiencequestion{這個一階導數的縮並表示式就是拉普拉斯算符嗎？}
\lecturerresponse{不是。拉普拉斯算符包含二階空間導數。這個表示式是由一階導數乘積構成的拉格朗日密度。\lecturetimestamp{01:24:38}}
\end{classroomqa}

標量場作用密度在物理事件處的值與參考系無關。適當的洛倫茲變換行列式絕對值為 1，因此積分測度\(d^4x\)也保持不變。於是，
\begin{equation}
S=\int d^4x\,\mathcal L
\end{equation}
也是不變量。與參考系無關的作用量，會導出在每個慣性參考系中形式相同的尤拉--拉格朗日方程。

\section{伽利略極限}

\begin{classroomqa}
\audiencequestion{牛頓力學的拉格朗日量也必須遵守相對性原理嗎？}
\lecturerresponse{是的，只是先前沒有特別強調。牛頓方程在伽利略變換下保持不變；伽利略變換是洛倫茲變換在\(c\to\infty\)下的極限。在這個極限中消失的是同時性的相對性。\lecturetimestamp{01:27:16}}
\end{classroomqa}

為了看出這個極限，先在沿\(x\)方向的推動中恢復\(c\)：
\begin{equation}
x'
=
\frac{x-vt}{\sqrt{1-v^2/c^2}},
\qquad
t'
=
\frac{t-vx/c^2}{\sqrt{1-v^2/c^2}}.
\label{eq:sr5-lorentz-c}
\end{equation}
令\(c\to\infty\)，可得
\begin{equation}
\boxed{
x'=x-vt,
\qquad
t'=t
}.
\label{eq:sr5-galilean}
\end{equation}
這就是伽利略變換。牛頓力學保留了相對性原理：定律在所有慣性參考系中形式相同；但它也假設所有觀察者共用一個絕對時間。

\begin{figure}[htbp]
\centering
\includegraphics[width=0.90\textwidth]{lecture_05_figure_22.png}
\caption{取\(c\to\infty\)前的洛倫茲推動，以及由此得到的伽利略變換\(x'=x-vt\)和\(t'=t\)。}
\lectureframe{01:29:40}
\end{figure}

\begin{classroomqa}
\audiencequestion{上標向量與下標向量的差異是否只有一個符號？}
\lecturerresponse{對此處採用的固定笛卡兒度規而言，對已知向量降指標只會改變時間分量的符號。不過，兩者的洛倫茲變換律是互為對偶的變換律，不只是任意的符號約定；度規負責將其中一種轉換成另一種。\lecturetimestamp{01:31:19}}
\end{classroomqa}

\begin{classroomqa}
\audiencequestion{每個看起來普通的座標函數都是標量嗎？}
\lecturerresponse{不是。在每個座標點上取一個數值還不夠；這個數值必須在每個參考系中，都代表同一事件處的同一物理量。能量密度就是一個反例：推動一個盒子會改變測得的能量，也會改變測得的體積，因此不同觀察者測得的密度不會相同。\lecturetimestamp{01:33:23}}
\end{classroomqa}

四維速度的一個分量也能說明這點。完整的四維速度是向量，但它的時間分量會隨推動而改變，因此本身不是標量。

\section{克萊因--戈登場}

\subsection{諧振子的場論類比}

回到標量場的作用密度，加入最簡單的二次勢：
\begin{equation}
\boxed{
\mathcal L
=
\frac12
\left[
(\partial_t\phi)^2
-|\boldsymbol{\nabla}\phi|^2
-\mu^2\phi^2
\right]
}.
\label{eq:sr5-kg-density}
\end{equation}
這與\begin{equation}
L_{\mathrm{osc}}
=
\frac12\dot q^2-\frac12\mu^2q^2.
\end{equation}的場論形式相似。每個空間點在局部都像一個振子，而梯度項則讓相鄰的點彼此耦合。

\begin{figure}[htbp]
\centering
\includegraphics[width=0.90\textwidth]{lecture_05_figure_23.png}
\caption{二次標量場作用密度及其諧振子原型。}
\lectureframe{01:38:20}
\end{figure}

場的尤拉--拉格朗日方程為
\begin{equation}
\partial_\mu
\left[
\frac{\partial\mathcal L}
{\partial(\partial_\mu\phi)}
\right]
-
\frac{\partial\mathcal L}{\partial\phi}
=0.
\label{eq:sr5-field-el}
\end{equation}
將它套用於\eqref{eq:sr5-kg-density}，得到
\begin{equation}
\boxed{
\partial_t^2\phi-\nabla^2\phi+\mu^2\phi=0
}.
\label{eq:sr5-kg}
\end{equation}
這就是克萊因--戈登方程。在時間分量為負、空間分量為正的號差慣例下，並採用\(\Box=\partial_\mu\partial^\mu=-\partial_t^2+\nabla^2\)，同一方程可寫成
\((\Box-\mu^2)\phi=0\)。

\begin{figure}[htbp]
\centering
\includegraphics[width=0.90\textwidth]{lecture_05_figure_24.png}
\caption{克萊因--戈登方程寫在導出它的作用密度下方。}
\lectureframe{01:40:00}
\end{figure}

從歷史上看，這個方程最初是作為相對論性單粒子波動方程來研究的。這種詮釋遇到了一些困難，尤其是難以得到正的機率密度。現代而言，它最適合放在相對論性場論中理解；它支配自由標量場，也描述具相互作用標量場的自由極限。

\subsection{平面波模態}

由於\eqref{eq:sr5-kg}是係數為常數的線性方程，可以設一個平面波試探解
\begin{equation}
\phi(x)
=
\exp\!\left[
-i\omega t+i\mathbf{k}\cdot\mathbf{x}
\right].
\label{eq:sr5-plane-wave}
\end{equation}
其中
\begin{equation}
k_\mu=(-\omega,k_x,k_y,k_z),
\end{equation}
指數為\(ik_\mu x^\mu\)。

\begin{figure}[htbp]
\centering
\includegraphics[width=0.90\textwidth]{lecture_05_figure_25.png}
\caption{以分量形式及\(e^{ik_\mu x^\mu}\)表示的平面波試探解。}
\lectureframe{01:44:00}
\end{figure}

微分後得到
\begin{equation}
\partial_t^2\phi=-\omega^2\phi,
\qquad
\nabla^2\phi=-|\mathbf{k}|^2\phi.
\end{equation}
將結果代入\eqref{eq:sr5-kg}，可得
\begin{equation}
\left(
-\omega^2+|\mathbf{k}|^2+\mu^2
\right)\phi=0.
\end{equation}
因此，非零模態必須滿足
\begin{equation}
\boxed{
\omega^2=|\mathbf{k}|^2+\mu^2
},
\qquad
\omega=\pm\sqrt{|\mathbf{k}|^2+\mu^2}.
\label{eq:sr5-dispersion}
\end{equation}

\begin{figure}[htbp]
\centering
\includegraphics[width=0.90\textwidth]{lecture_05_figure_26.png}
\caption{代入平面波試探解，得到相對論性色散關係。}
\lectureframe{01:47:00}
\end{figure}

在自然單位制中，\(\mu\)扮演質量的角色。對量子激發恢復常數後，
\begin{equation}
E=\hbar\omega,
\qquad
\mathbf{p}=\hbar\mathbf{k},
\qquad
\mu=\frac{mc}{\hbar},
\end{equation}
便可將\eqref{eq:sr5-dispersion}改寫為
\begin{equation}
E^2=p^2c^2+m^2c^4.
\end{equation}
正頻率支與負頻率支都屬於經典解空間。一般解是這些模態的疊加。

\begin{classroomqa}
\audiencequestion{克萊因--戈登方程是如何得到的？}
\lecturerresponse{它直接由二次洛倫茲標量密度出發，套用場的尤拉--拉格朗日方程而得。選這個例子，是因為這種簡單的二次作用密度會導出線性場方程。\lecturetimestamp{01:47:59}}
\end{classroomqa}

\begin{classroomqa}
\audiencequestion{自然界中有哪些東西遵循克萊因--戈登方程？}
\lecturerresponse{自由的經典標量場遵循此方程。希格斯場的漲落是標量，在自由近似下會化為克萊因--戈登形式；不過，實際的希格斯場還包含相互作用與規範耦合。\lecturetimestamp{01:50:41}}
\end{classroomqa}

\section{標量背景造成的質量位移}

現在回到粒子作用量：
\begin{equation}
S_{\mathrm p}
=
-\int d\tau\,
\bigl[m+g\phi(x)\bigr].
\label{eq:sr5-final-particle-action}
\end{equation}
如果\(\phi\)具有常數背景值\(\phi_0\)，粒子的運動就如同它的質量是
\begin{equation}
\boxed{
m_{\mathrm{eff}}=m+g\phi_0
}.
\label{eq:sr5-effective-mass}
\end{equation}
若裸質量項\(m\)為零，標量背景便能提供乘在固有時前的全部係數。

\begin{figure}[htbp]
\centering
\includegraphics[width=0.90\textwidth]{lecture_05_figure_27.png}
\caption{上方依賴場值的粒子作用量，以及下方簡潔的場尤拉--拉格朗日方程，將本講的兩個部分連結起來。}
\lectureframe{01:51:30}
\end{figure}

這是刻意簡化的希格斯質量生成模型。在標準模型中，費米子質量來自與希格斯真空期望值的湯川耦合，示意性地可寫成 \(m_f\propto y_f v\)；\(W\) 和 \(Z\) 的質量則來自規範協變的希格斯動能項。這個玩具模型中的世界線耦合掌握了核心概念--非零的純量背景會改變慣性質量--但並未重現完整的電弱理論。先前建立的互易關係仍然成立：對同一個相互作用項關於場作變分，粒子便成為場的源。場作用於粒子，粒子也會反過來作用於場。

\subsection{複指數為何有用}

場在物理上不必是複數。尤拉公式給出
\begin{equation}
e^{i(\omega t-\mathbf{k}\cdot\mathbf{x})}
=
\cos(\omega t-\mathbf{k}\cdot\mathbf{x})
+i\sin(\omega t-\mathbf{k}\cdot\mathbf{x}).
\end{equation}
由於方程是實係數的線性方程，實部和虛部分別都是實解。它們是相位相差四分之一週期的波，而任意實係數線性組合仍是解。

\begin{classroomqa}
\audiencequestion{取實部時，虛部會以某種方式相加而變成零嗎？}
\lecturerresponse{不會。複指數將兩個獨立的實解，也就是餘弦和正弦，合在一起。可以取其中任一部分，也可以組成實線性組合。使用指數形式，是因為微分運算會變成乘法。\lecturetimestamp{01:53:29}}
\end{classroomqa}

\begin{figure}[htbp]
\centering
\includegraphics[width=0.90\textwidth]{lecture_05_figure_28.png}
\caption{將複平面波拆成餘弦波與正弦波，圖中以兩個相位不同的實波呈現。}
\lectureframe{01:56:00}
\end{figure}

\begin{classroomqa}
\audiencequestion{超導體中的電磁場獲得質量，是同一類現象嗎？}
\lecturerresponse{兩者密切相關：在超導體中，凝聚態會使電磁場在介質中具有有效質量，因而產生邁斯納效應。完整說明需要引入電磁場，這部分留待後續討論。\lecturetimestamp{01:56:14}}
\end{classroomqa}

\begin{classroomqa}
\audiencequestion{能否把希格斯機制想像成希格斯粒子在空間中飛來飛去、把質量分給其他粒子？}
\lecturerresponse{不能。關鍵是希格斯場的背景值，而不是一股由個別希格斯粒子組成的流。在這個玩具模型的作用量中，背景值只會改變那個作為粒子質量的係數。\lecturetimestamp{01:57:31}}
\end{classroomqa}

\begin{classroomqa}
\audiencequestion{同一個純量背景透過不同耦合，會產生不同質量嗎？}
\lecturerresponse{會。質量的改變正比於耦合常數 \(g\)。在標準模型中，不同的湯川耦合正是不同費米子從同一個希格斯背景獲得不同質量的原因。\lecturetimestamp{01:58:04}}
\end{classroomqa}

\begin{classroomqa}
\audiencequestion{\(\phi\)本身就是質量項嗎？}
\lecturerresponse{在這個玩具模型中，\(g\phi_0\)就所有力學作用而言，都相當於質量係數的一項貢獻。耦合常數和背景值都必須納入考量。\lecturetimestamp{01:58:42}}
\end{classroomqa}

\begin{classroomqa}
\audiencequestion{如果希格斯背景消失，所有粒子都會變成無質量嗎？}
\lecturerresponse{不會。電子、夸克和弱規範玻色子會失去通常由電弱作用提供的質量；中微子的情況則取決於賦予中微子質量的機制。即使希格斯場的背景值消失，純量部門本身也不一定會變成無質量。質子和中子的質量大多來自 QCD 的束縛與禁閉效應，因此並非直接由希格斯場生成。\lecturetimestamp{01:58:51}}
\end{classroomqa}

然而，對原子的影響仍然十分劇烈。波耳尺度的關係為
\begin{equation}
a_0\sim\frac{1}{\alpha m_e}
\end{equation}
這裡採用自然單位制。若電子質量趨近於零，通常的原子長度尺度便會發散；一般的局域原子將不復存在。質子和中子受影響的程度則小得多，因為它們的質量主要來自 QCD 的動力學能量。

至此，這些內容串成了一條完整的論證。洛倫茲幾何決定觀測到的波長。一個作用量中的一個混合項，便能使兩個動力學系統彼此影響。純量場既是粒子的勢，也能成為粒子的源，進而形成場。度規與指標記號讓洛倫茲不變性一目了然，而最簡單的二次不變量密度會產生克萊因--戈登波。最後，常數純量背景會透過同一個耦合改變質量。貫穿全文、也將繼續採用的方法是：構造不變作用量，對每個動力學變數作變分，再從所得方程讀出其中的相互作用物理。